\documentclass[12pt, a4paper]{article}
\usepackage[utf8]{inputenc}
\usepackage{graphicx}
\usepackage{geometry}
\usepackage{hyperref}
\usepackage{bookmark}
\usepackage{float}
\usepackage{titlesec}
\usepackage{xcolor}
\usepackage{setspace}
\usepackage{amsmath}
\usepackage{listings}

% Page Geometry setup
\geometry{
    a4paper,
    total={170mm,257mm},
    left=25mm,
    top=25mm,
    right=25mm,
    bottom=25mm,
}

% Code Listing Configuration
\definecolor{codegreen}{rgb}{0,0.6,0}
\definecolor{codegray}{rgb}{0.5,0.5,0.5}
\definecolor{codepurple}{rgb}{0.58,0,0.82}
\definecolor{backcolour}{rgb}{0.95,0.95,0.92}

\lstdefinestyle{mystyle}{
    backgroundcolor=\color{backcolour},   
    commentstyle=\color{codegreen},
    keywordstyle=\color{magenta},
    numberstyle=\tiny\color{codegray},
    stringstyle=\color{codepurple},
    basicstyle=\ttfamily\footnotesize,
    breakatwhitespace=false,         
    breaklines=true,                 
    captionpos=b,                    
    keepspaces=true,                 
    numbers=left,                    
    numbersep=5pt,                  
    showspaces=false,                
    showstringspaces=false,
    showtabs=false,                  
    tabsize=2
}
\lstset{style=mystyle}

% Set line spacing to 1.5 for formal report standard
\onehalfspacing

% Title Page Information
\title{\textbf{\Huge UIDAI Hackathon Project Report} \\ \vspace{0.5cm} \Large Nexus Analytics: Strategic Intelligence Suite}
\author{\textbf{Team Name} \\ [0.5cm] \textit{Submitted for UIDAI Hackathon}}
\date{\today}

\begin{document}

\maketitle
\thispagestyle{empty}

\begin{abstract}
    \noindent \textbf{Nexus Analytics} implements a high-precision, AI-driven framework for analyzing large-scale Aadhar enrolment datasets. This report details our deep-dive analysis into demographic shifts, operational anomalies, and growth patterns. Moving beyond descriptive statistics, we identify critical "Velocity Fraud" vectors and "Saturation" points, providing a data-backed roadmap for optimizing the next phase of the enrolment ecosystem.
    
    \vspace{0.5cm}
    \noindent \textbf{Live Dashboard:} \url{http://127.0.0.1:5000/}
\end{abstract}

\newpage
\tableofcontents
\newpage

% ---------------------------------------------------------
% SECTION 1: Problem Statement and Approach
% ---------------------------------------------------------
\section{Problem Statement and Approach}

\subsection{Problem Statement}
The Unique Identification Authority of India (UIDAI) manages the world's largest biometric ID system. As the system matures, administrators require deep, actionable insights rather than simple volume metrics. The core problem addressed is the "Operational Blindspot" — the difficulty in distinguishing between:
\begin{itemize}
    \item \textbf{High Efficiency vs. Fraud}: Is a high-volume center efficient, or is it using injection scripts?
    \item \textbf{Growth vs. Saturation}: Is a district's low growth a failure of effort, or simply a result of reaching 100\% adult saturation?
    \item \textbf{Missing Demographics}: Which specific pockets are failing to capture birth rates?
\end{itemize}

\subsection{Proposed Technical Approach}
To address this, we developed **Nexus Analytics**, an AI-driven Intelligence Engine. Our approach centers on a "4-Layer Deep Analysis" framework:
\begin{enumerate}
    \item \textbf{Statistical Normalization}: We apply Z-Score standardization to district-level data to enable fair comparisons between large urban centers and small rural octs.
    \item \textbf{Pattern Recognition}: We correlate "Child Ratio" with "Total Volume" to identify "Child Focus Areas" versus "Missed Opportunities."
    \item \textbf{Predictive Modeling}: Using Linear Regression on time-series data to forecast future inventory and resource needs.
    \item \textbf{Actionable Intelligence}: Converting raw stats into English-language "Executive Advice" (e.g., "Shift budget to Child Campaigns").
\end{enumerate}

\begin{figure}[H]
    \centering
    % Architecture Diagram Image
    \includegraphics[width=0.7\textwidth]{arch.png}
    \caption{System Approach: Converting Raw CSV Data into Strategic Action.}
    \label{fig:arch}
\end{figure}

% ---------------------------------------------------------
% SECTION 2: Datasets Used
% ---------------------------------------------------------
\newpage
\section{Datasets Used}

For this analysis, we utilized the official Aadhar enrolment and update datasets provided for the hackathon.

\subsection{Data Source}
The system ingests the provided CSV shards:
\begin{itemize}
    \item \texttt{api\_data\_aadhar\_enrolment\_0\_500000.csv}
    \item \texttt{api\_data\_aadhar\_enrolment\_500000\_1000000.csv}
    \item \texttt{api\_data\_aadhar\_enrolment\_1000000\_1006029.csv}
\end{itemize}

\subsection{Key Attributes}
We specifically focused on the following attributes for our analytical models:
\begin{itemize}
    \item \textbf{Temporal}: \texttt{date} (Used for Time-Series Analysis and Growth Momentum).
    \item \textbf{Geospatial}: \texttt{state}, \texttt{district} (Used for Regional grouping and Heatmaps).
    \item \textbf{Demographic Buckets}:
    \begin{itemize}
        \item \texttt{age\_0\_5}: New birth/child enrolments (Key metric for future growth).
        \item \texttt{age\_5\_17}: School-age updates (Mandal biometrics).
        \item \texttt{age\_18\_greater}: Adult population (Used to calculate Saturation).
    \end{itemize}
    \item \textbf{Volume}: \texttt{total} (Aggregated enrolment counts).
\end{itemize}

% ---------------------------------------------------------
% SECTION 3: Methodology
% ---------------------------------------------------------
\section{Methodology}

\subsection{Data Cleaning and Preprocessing}
The raw data required significant preprocessing to ensure statistical validity:
\begin{enumerate}
    \item \textbf{Unified Time-Series}: The fragmented CSV files were concatenated into a single Pandas DataFrame. The \texttt{date} column was parsed to datetime objects.
    \item \textbf{Seasonal Filtering}: We filtered out data from specific start-up months (Jan/Feb) to remove initialization noise and focus on the "Active Growth" phase.
    \item \textbf{Entity Normalization}: State and District names were stripped of whitespace and normalized to ensure correct aggregation.
    \item \textbf{Null Imputation}: Numeric columns contained \texttt{NaN} values which were imputed with 0 to prevent propogation errors in summation.
\end{enumerate}

\subsection{Transformations and Algorithms}
\subsubsection{1. Z-Score for Velocity Fraud}
To detect anomalies, we calculate the Z-Score for daily district volumes.
\[ Z = \frac{x - \mu}{\sigma} \]
Districts with $|Z| > 2.5$ are flagged as "Statistical Anomalies," suggesting potential batch-injection fraud.

\subsubsection{2. Growth Momentum Calculation}
We calculate the Month-over-Month (MoM) growth momentum:
\[ M = \frac{V_{current} - V_{previous}}{V_{previous}} \times 100 \]
This transforms absolute volume into a "Velocity" vector, allowing us to classify states as "Accelerators" or "Decelerators."

\subsubsection{3. Predictive Forecasting}
We utilize a Least Squares Polynomial Fit (Degree 1) on the monthly aggregated time-series to project the trend line $y = mx + c$ for the upcoming quarter.

% ---------------------------------------------------------
% SECTION 4: Data Analysis and Visualisation
% ---------------------------------------------------------
\newpage
\section{Data Analysis and Visualisation}

\subsection{Key Findings}
\begin{itemize}
    \item \textbf{Adult Saturation}: Top-tier states have reached a "Saturation Point" where adult updates constitute $<30\%$ of total volume.
    \item \textbf{Velocity Fraud}: We detected specific districts reporting $>1000$ enrolments/day, which is statistically impossible for a single center, indicating rigorous need for "Speed Limit" protocols.
    \item \textbf{Child Laggards}: Despite high overall volumes, certain districts show a \texttt{child\_ratio} of $<5\%$, indicating a failure to capture new births.
\end{itemize}

\subsection{Visualisations}
(Note: All 11 dashboard visualizations are included below).

\subsubsection{1. Demographic Breakdown (Pie Chart)}
A high-level view of the overall distribution, highlighting the proportion of child vs. adult enrolments.
\begin{figure}[H]
    \centering
    \includegraphics[width=0.7\textwidth]{pie.png} 
    \caption{Overall Demographic Distribution (0-5, 5-17, 18+)}
\end{figure}

\subsubsection{2. Top State Performance (Bar Chart)}
Ranking the top 10 states by total volume to identify key contributors.
\begin{figure}[H]
    \centering
    \includegraphics[width=0.7\textwidth]{states.png}
    \caption{Top 10 States by Enrolment Volume}
\end{figure}

\subsubsection{3. Historical Volume Trend (Area Chart)}
Overall monthly enrolment volume since dataset inception.
\begin{figure}[H]
    \centering
    \includegraphics[width=0.7\textwidth]{trend.png}
    \caption{Historical Enrolment Trend}
\end{figure}

\subsubsection{4. Monthly Age Trends (Stack Bar)}
Monthly breakdown showing which age cohort is driving specific seasonal spikes.
\begin{figure}[H]
    \centering
    \includegraphics[width=0.7\textwidth]{agetrend.png}
    \caption{Monthly Trend by Age Group}
\end{figure}

\subsubsection{5. Growth Momentum Analysis}
Identifies which states are gaining speed versus those losing momentum.
\begin{figure}[H]
    \centering
    \includegraphics[width=0.7\textwidth]{momentum.png}
    \caption{State Momentum: Accelerators vs Decelerators}
\end{figure}

\subsubsection{6. Demographic Saturation (Heatmap)}
Visualizes the shift from Adult-heavy to Child-heavy enrolment in mature states.
\begin{figure}[H]
    \centering
    \includegraphics[width=0.7\textwidth]{heatmap.png}
    \caption{Demographic Breakdown by Age Group across Sectors}
\end{figure}

\subsubsection{7. Child Enrolment Hotspots (Bar Chart)}
Identifying the specific districts leading in 0-5 child enrolment capture.
\begin{figure}[H]
    \centering
    \includegraphics[width=0.7\textwidth]{hotspots.png}
    \caption{Top 10 Districts: Child Enrolment Hotspots}
\end{figure}

\subsubsection{8. Efficiency Matrix (Scatter Plot)}
Correlates volume ($x$) with child ratio ($y$). Top-left outliers are "Stars" (efficient child capture).
\begin{figure}[H]
    \centering
    \includegraphics[width=0.7\textwidth]{patterns.png}
    \caption{AI Pattern Analysis: Efficiency Matrix}
\end{figure}

\subsubsection{9. Deep Trend Velocity (Line Chart)}
Comparitive speed of growth between Child vs Adult segments.
\begin{figure}[H]
    \centering
    \includegraphics[width=0.7\textwidth]{deeptrend.png}
    \caption{Deep Trend Velocity Analysis}
\end{figure}

\subsubsection{10. Anomaly Detection (Velocity Fraud)}
Red bubbles indicate districts exceeding the statistical threshold ($Z > 2.5$).
\begin{figure}[H]
    \centering
    \includegraphics[width=0.7\textwidth]{anomaly.png}
    \caption{Anomaly Detection: Volume Spikes}
\end{figure}

\subsubsection{11. Predictive Forecasting}
AI-Projected load for the next 3 months.
\begin{figure}[H]
    \centering
    \includegraphics[width=0.7\textwidth]{forecast.png}
    \caption{3-Month Enrolment Forecast Trend}
\end{figure}

% ---------------------------------------------------------
% APPENDIX: Code Files
% ---------------------------------------------------------
\newpage
\appendix
\section{Code Files}
As per the requirement, the source code used for this analysis is included below.

\subsection{Analytical Engine (app.py)}
\lstinputlisting[language=Python, caption=Main Analytical Engine]{app.py}

\newpage
\subsection{Visualisation Logic (script.js)}
\lstinputlisting[language=JavaScript, caption=Frontend Charting Logic]{script.js}

\end{document}
